\chapter{Quick Start Guide}
\label{ch:quick_start}

Follow these steps to set up a basic ducking configuration in your DAW.

\section{Basic Setup Example (Kick \& Bass)}
Goal: Duck the bass track whenever the kick drum hits.

1. \textbf{Insert Duq:} Place Duq as an insert effect on your \textbf{Bass} track.
2. \textbf{Route MIDI:}
   \begin{itemize}
       \item Create a new MIDI track or use your existing Kick MIDI track.
       \item Route the MIDI output of the Kick track to the Bass track where Duq is inserted.
       \item Ensure Duq is receiving MIDI notes (indicated by the UI trigger feedback).
   \end{itemize}
3. \textbf{Select Envelope:} In the Envelope List (left sidebar), select an envelope or create a new one.
4. \textbf{Assign Note:} Click the MIDI note display on the envelope row and press the MIDI note corresponding to your kick drum (e.g., C3).
5. \textbf{Shape the Curve:}
   \begin{itemize}
       \item Double-click the grid to add points.
       \item Drag points to define the ducking shape (usually starting at 1.0, dipping to 0.0, and returning to 1.0).
       \item Adjust the curve tension by dragging the anchors between points.
   \end{itemize}
6. \textbf{Adjust Parameters:}
   \begin{itemize}
       \item Set \textbf{Rate} to sync or a specific Hz.
       \item Increase \textbf{Depth} to 100\% for total silence at the dip.
       \item Adjust \textbf{Lookahead} if the ducking feels slightly late (common for transient preservation).
   \end{itemize}
7. \textbf{Verify:} Play back your project and observe the gain reduction meter in Duq.
